% Lösungen Einheit 2 von 5 – Bernoulli-Formel (zusammenbau v0.1)
\einheitenkopf{Einheit 2 von 5 · Bernoulli-Formel}

\erg{13}{a) $P(X = 4)$ – „genau“ verlangt die Einzelwahrscheinlichkeit \quad b) $(6 \text{ über } 2) \cdot 0{,}35^{2} \cdot 0{,}65^{4} \approx 0{,}3280$ \quad c) $(8 \text{ über } 3) = 56$; $P = 56 \cdot 0{,}12^{3} \cdot 0{,}88^{5} \approx 0{,}0511$ \quad d) $0{,}92^{30} \approx 0{,}0820$ \quad e) Treffer „nicht erscheinen“ mit $0{,}10$, genau 2: $(20 \text{ über } 2) \cdot 0{,}10^{2} \cdot 0{,}90^{18} \approx 0{,}2852$ \quad f) $\left(\frac{5}{6}\right)^4 + 4 \cdot \frac{1}{6} \cdot \left(\frac{5}{6}\right)^3 = \frac{1\,125}{1\,296} \approx 0{,}8681$ \quad g) $k = 4$, $n = 7$, $1 - p = \frac{3}{5}$: $P(X = 4) = (7 \text{ über } 4) \cdot \left(\frac{3}{5}\right)^3 \cdot \left(\frac{2}{5}\right)^4$ \quad h) $P(X = 40) \approx 0{,}0716$ (Rechner: n = 120, p = 0{,}35, k = 40; X: Anzahl) \quad i) Trefferwahrscheinlichkeit $\frac{7}{20} = 0{,}35$; $P(X = 7) \approx 0{,}1844$ \quad j) $P(X = 6) = P(X \le 6) - P(X \le 5) \approx 0{,}1712$; $P(X < 6) = P(X \le 5) \approx 0{,}4353$}
\erg{14}{a) Fehler: der Binomialkoeffizient fehlt, Jonas rechnet nur einen Pfad. Richtig: $(5 \text{ über } 2) \cdot \left(\frac{1}{6}\right)^2 \cdot \left(\frac{5}{6}\right)^3 = 10 \cdot 0{,}0161 \approx 0{,}1608$ \quad b) Jeder Pfad ist ein Produkt aus zwei Faktoren p und einem Faktor 1 − p; die Reihenfolge der Faktoren ändert das Produkt nicht. \quad c) $0{,}96^{52} + 52 \cdot 0{,}04 \cdot 0{,}96^{51} \approx 0{,}3791$ \quad d) Trefferwahrscheinlichkeit 0,15, etwa ein Glücksrad mit Gold; sechs Drehungen, genau zweimal Gold.}
